\clearpage
\section{Выводы}

В ходе лабораторной работы построена корректная геометрическая модель усечённого правильного тетраэдра и реализовано её отображение на двумерной плоскости. Рассмотрены различия перспективной и ортографической проекций, поворот пространственных вершин, определение видимых граней и преобразование координат в экранную систему.

Получены практические навыки работы со стартовым проектом Vulkan, обработки окна и ввода через GLFW и создания интерактивного интерфейса с Dear ImGui. Проверены правильность граней, замкнутость поверхности и сохранение геометрии при повороте. Освоены запись буферов команд и последовательное выполнение проходов рисования.

Базовый вариант № 4 выполнен. Переключение проекций, вращение и масштаб отображения доступны в реальном времени. Выбранный способ визуализации достаточен для одного выпуклого объекта. Для дальнейшего расширения можно добавить трансформации положения и растяжения, управление цветом, сложную анимацию и вывод нескольких моделей с отдельными ресурсами.
